% !TEX root = ../main.tex

\chapter{Introduction}
\minitoc

Bulk Synchronous Parallel (BSP) execution underpins many of the most expensive and consequential
computational workloads in use today, from large-scale scientific simulations to distributed
training of frontier machine learning models.
These workloads run for weeks to months on thousands to tens of thousands of accelerators, so even
small efficiency losses translate into major costs in energy, compute time, and delayed
time-to-discovery.

The same synchronized execution model that makes BSP effective also creates two persistent
efficiency pressure points.
On the data path, coordinated checkpoints and analysis outputs generate bursty I/O that strains
storage and slows downstream analysis.
On the compute path, global synchronization makes end-to-end progress sensitive to stragglers and
runtime anomalies, and also amplifies correctness risks when failures or corrupted state propagate
across ranks.


\begin{tcolorbox}[thesisbox]
  A large class of efficiency bottlenecks in BSP systems arise from conditions that are unknowable
  before runtime.
  Through two studies spanning the data and compute paths, this thesis demonstrates that such
  conditions are systematically addressable through adaptive mechanisms if the requisite feedback
  loops exist, and perpetuate if they do not.
  We present ORCA as general-purpose infrastructure for closing these loops at scale.
\end{tcolorbox}

\section{The Problem of Runtime Uncertainty}

We investigated BSP efficiency along two paths: data-path indexing for range queries, and
compute-path placement for adaptive mesh codes.
Both are shaped by runtime uncertainty---conditions unknowable before execution.
CARP builds an online mechanism to exploit this, while the AMR study highlights a broader class of
runtime uncertainties, the tooling bottlenecks that impede their resolution, and the need for a
general-purpose solution.

\subsection{CARP: Range Query-Optimized Indexing for Streaming Data}

CARP (\emph{Continuous Adaptive Range Partitioner}) is an in-situ partitioner for efficient range
queries on scientific data, e.g.\ all particles within a given energy range.
In-situ range partitioning is challenging because scientific key distributions are highly skewed,
evolve over time, and make any static partitioning load-imbalanced.
CARP uses a reconstructed view of key distributions at runtime to drive load-balanced partitioning.
An adaptive mechanism tracks key distributions as they evolve to compute new partition boundaries,
without reassigning existing data.
As a result, CARP produces a partially ordered layout in a single streaming pass, deferring some
compaction to query time.
In evaluation, CARP approximates the query performance of a full sort without any ingestion
overhead, making it up to $4.9\times$ faster than post-processing approaches.

\subsection{Lessons from Profiling and Optimizing Placement in AMR Codes}

On the compute path, we studied placement in \emph{adaptive mesh refinement} (AMR) codes, a class
of highly dynamic simulations where work characteristics evolve to track physical phenomena.
Poor utilization in these codes is widely attributed to load imbalance.
We sought to develop telemetry-driven placement policies that would adapt to evolving workload
patterns, but found genuine placement effects routinely confounded by unrelated artifacts:
fail-slow hardware, fabric pathologies, and tuning issues.
Diagnosis consumed months of iterative debugging even though mitigations, once identified,
typically took days.
With reliable telemetry established, we found that placement exposed a fundamental tradeoff between
compute load balance and communication locality, and designed CPLX, a tunable placement policy that
achieves up to 21.6\% runtime reduction over tuned baselines across 512--4096 ranks.

\subsection{The Feedback Gap}

The two studies illustrate two modes of dealing with runtime uncertainty.
In CARP, an automated policy tracks attribute distributions and adapts partitioning in real
time---effective, but a fixed-function mechanism for a single use case.
In the AMR study, a human operator drives diagnosis and intervention, but the iteration latency of
existing tools turns what should be a rapid cycle into one spanning months.
Each artifact took months to isolate, while mitigations once identified typically took days.
The core problem was structural:

\begin{enumerate}
  \item \textbf{Every telemetry mode requires re-execution.}
        Profiling summaries are too coarse, traces too voluminous, and custom instrumentation requires
        recompilation---each candidate hypothesis costs a full application run.

  \item \textbf{Trace ETL delays insight.}
        Even when traces are collected, their volume creates an ingestion and analysis bottleneck that can
        take longer than the run itself.

  \item \textbf{Post-hoc insight delays intervention.}
        There is no way to know that an anomaly has occurred until after execution completes and traces are
        analyzed.
        For workloads that run for weeks, problems go undetected for the duration of the run.
\end{enumerate}

The practical consequence is that rigorous performance analysis is often not undertaken at all.
When traces are collected, they become archival datasets analyzed by groups without access to the
runtime environment that produced them.
These limitations are not unique to our work---the same persistent class of structural bottlenecks
recurs across two decades of HPC performance literature.

\section{Solution: View-Driven and Steerable Observability}

Addressing BSP efficiency requires infrastructure that can close feedback loops in real time with
\emph{view-driven} and \emph{steerable} observability.

\paragraph{View-driven}
The unit of feedback should be a materialized view, not raw data.
Tracing captures comprehensive event logs and defers analysis to post-processing, creating a cycle
where more data begets more ETL work---cost scales with data volume regardless of what the operator
actually needs to know.
The alternative is to invert this relationship: specify a hypothesis, and materialize only the view
that answers it, via in-situ reduction over distributed state.
Cost then scales with the question, not the data.
CARP already operates this way---it materializes a view of the global key distribution to drive
partitioning, not a dataset to be analyzed later.
The AMR study needed exactly this capability but lacked it, and paid for it in months of iteration
latency.

\paragraph{Steerable}
A single observability substrate should span a continuous spectrum from lightweight summaries
sufficient for routine monitoring to detailed per-rank per-timestep telemetry for deep diagnosis,
with transitions between them available at runtime without re-execution or retooling.
This replaces the current fragmented choice between profiling, tracing, and custom instrumentation.
Steerability extends beyond observation to execution itself: views inform interventions---adjusting
placement, tuning parameters, modifying partitioning---and the impact of those interventions must
be assessed through further observation.
Who drives this process is orthogonal to the infrastructure: it may be a human operator, an
automated policy, or a combination of both.

\subsection{Status Quo and Limitations}

\paragraph{Tracing and post-hoc analytics}
In HPC, tools like TAU and Score-P must be configured before execution and collect fine-grained
traces, with analysis happening separately through post-processing.
The PyTorch ecosystem comes closest to a complete feedback loop, with Kineto for per-node trace
collection, Dynolog as a control plane, and HTA for trace analytics.
But traces are voluminous per-rank JSON logs with no awareness of synchronization structure,
requiring format-specific ETL before any analysis can begin.
The on-disk layout offers no structure to support selective access: per-rank logs must be parsed in
full for any query, with working sets scaling with ranks $\times$ timesteps.
Dynolog has no notion of timestep-coordinated collection, HTA is a single-node Pandas-based tool
that does not scale to cluster-level analysis, and collection is on-demand rather than always-on.
In both cases, the feedback gap is driven by three compounding costs: intrinsic trace volume,
format-dictated ETL, and suboptimal on-disk layouts.

\paragraph{Cloud observability}
Distributed systems monitoring offers an alternative lineage: Prometheus for metrics, Kafka and
Flink for event processing.
While viable and deployed for ML training workloads, these tools have fundamental limitations for
tightly-coupled parallel applications.
They do not preserve BSP causal structure---even the simplest problem of identifying a straggler at
a collective requires knowing when all ranks' telemetry is available, and then joining it by that
collective, both challenging for event-based analytics.
They cannot leverage HPC fabrics for in-situ reduction, leading to unnecessary data movement and
collection overhead.
And they do not enable steerability, as they provide no control semantics with meaningful
consistency guarantees.

\subsection{Why Now}

\paragraph{Dataframe-Style Trace Analytics}
Performance analysis increasingly looks like dataframe-shaped analytics---columnar formats, SQL
queries, aggregations over structured data.
Better formats such as Parquet, and better on-disk layouts informed by query patterns, enable
selective access and dramatically faster analytics: even offline, columnar trace analytics can be
orders of magnitude faster than existing tools.
In-situ reduction compounds this further by eliminating the intermediate dataset entirely.

\paragraph{Columnar Ecosystem Maturity}
Columnar formats such as Arrow and Parquet provide efficient in-memory and on-disk representations
for structured telemetry.
Embeddable SQL engines like DataFusion and DuckDB bring analytical query capabilities without
external infrastructure.
Fabric-aware RPC libraries such as Mercury over libfabric enable writing portable high-performance
transport code across HPC interconnects.

\paragraph{ML-Driven Demand}
Large-scale training runs are driving bespoke observability pipelines for specific global
views---correctness checks, straggler detection, throughput monitoring.
Each is a custom solution to what should be a general capability.
A single large training run justifies the investment; the pattern points toward a unified
substrate.

\section{ORCA: Real-time Observability and Control for BSP Workloads}

ORCA (\emph{Observability with Real-time Control and Aggregation}) is a tree-based overlay network
(TBON) for view-driven and steerable BSP observability, comprising a controller and a tier of
dedicated aggregators attached to a parallel application.
A SQL-based in-situ analytics interface allows telemetry, originating as columnar streams, to be
transformed en route to configurable sinks.
An asynchronous control plane provides atomic and timestep-consistent control semantics.
A key contribution is the \emph{timestep dataframe} abstraction, comprising all ranks' telemetry
for a single timestep.
Because global synchronization already divides BSP execution into causally independent windows,
this is the natural unit of in-situ analysis, and maps directly to \emph{(timestep,
  rank)}-partitioned Parquet for offline access.
Queries are compiled and plan operators pushed down the TBON tiers to minimize data movement.
On a real AMR code at 4096 ranks, ORCA collects detailed traces at 1–2\% overhead versus 56–81\%
for existing tracers, executes queries orders of magnitude faster, and reduces data volume by
multiple orders of magnitude through in-situ analytics.
In an agentic evaluation of the closed-loop workflow, a performance anomaly that took months to
diagnose manually is localized in 27 minutes.

\subsection{List of Contributions}

\paragraph{CARP}
\begin{itemize}
  \item An adaptive range partitioner for unknown and highly skewed key distributions.
        CARP ingests data in a single streaming pass and is $2.8$--$4.9\times$ faster than post-processing
        approaches.

  \item Query performance comparable to fully sorted layouts and up to two orders of magnitude faster than auxiliary bitmap indices.

  \item Shows that partially ordered layouts can be generated without reorganization, achieving perfect write amplification while enabling query performance comparable to fully sorted layouts.

  \item Together with DeltaFS's hash partitioning, completes the set of in-situ partitioning strategies for scientific data.
\end{itemize}

\paragraph{AMR Placement}
\begin{itemize}
  \item An end-to-end empirical study integrating collection, tuning, and placement, showing how artifacts such as overheating, fabric bugs, and operation ordering obscure placement effects.

  \item CPLX, a hybrid placement policy providing tunable control over the tradeoff between compute load balance and communication locality, achieving up to 21.6\% runtime reduction over tuned baselines on real AMR codes across 512--4096 ranks.

  \item Broader lessons on telemetry-driven performance analysis and the empirical, context-specific nature of performance in large-scale parallel systems, including a critical path framework for reasoning about stragglers in AMR codes.
\end{itemize}

\paragraph{ORCA}
\begin{itemize}
  \item The timestep dataframe abstraction, building on the insight that global synchronization divides BSP telemetry into causally independent windows.
        Timestep dataframes are persisted as \emph{(timestep, rank)}-partitioned Parquet, enabling up to
        $44.3\times$ smaller traces and up to $6{,}800\times$ faster post-mortem analysis.

  \item An incast-optimized RDMA transport for tree topologies, enabling detailed trace collection with under 2\% overhead versus 56--81\% for conventional tracers.

  \item A SQL-based programming interface with operator pushdowns, enabling 98.5--99.999\% data reduction at source for needle-in-the-haystack queries.

  \item A timestep-consistent control plane via a protocol called TS2PC (\emph{timestep-linked two-phase commit}).
        In an evaluation of closed-loop capabilities using an agentic debugging workflow, an LLM diagnosed
        a performance anomaly in 26 minutes that took months to isolate manually.
\end{itemize}

\section{Document Structure}

The remainder of this thesis is organized as follows.
\begin{itemize}
  \item \cref{chap:bg} provides background on BSP environments and efficiency bottlenecks.
  \item \cref{chap:carp} presents CARP and its evaluation on ingestion/query performance.
  \item \cref{chap:amr} presents the AMR diagnostic workflow and CPLX placement policy.
  \item \cref{chap:orca} presents ORCA and its observability/control design.
  \item \cref{chap:timeline} outlines remaining milestones and completion plan.
\end{itemize}
